%% ICAIF 2026 submission. ACM sigconf, anonymous, 8-page hard limit.
%% This is a self-contained compression of Paper_v8 to the conference format.
%% No supplementary material; all content fits in the 8-page body.

\documentclass[sigconf,anonymous,nonacm=true]{acmart}

%% Override the ACM acmart default section numbering off-switch so that
%% section headers carry numbers as in the conference template examples.
\settopmatter{printacmref=false}
\setcopyright{none}
\renewcommand\footnotetextcopyrightpermission[1]{}
\pagestyle{plain}

%% acmart already provides amsmath, amssymb, booktabs, multirow.
%% Redeclaring amssymb conflicts with acmart's \Bbbk definition.

%% Title and anonymized authors.
\title{A Continuous Hidden Markov Model for Synthetic Equity Returns:
  Emission-Family Ablation and Rank-Reordering Copula Extensions}

\author{Anonymous Author(s)}
\affiliation{%
  \institution{Submission to ICAIF 2026 (double-blind)}
  \city{}
  \country{}
}

\begin{document}

\begin{abstract}
We train a continuous hidden Markov model (CHMM) by Baum-Welch /
Expectation-Conditional-Maximization on raw annualized excess log returns and
show that, at a single moderate state resolution ($K = 18$), it reproduces all
three canonical stylized facts of equity returns (heavy tails, negligible
linear autocorrelation, persistent volatility clustering) without any jump
mechanism or return-space discretization. Three emission families are
compared under identical training: Gaussian (CHMM-N), Student-t with
per-state degrees of freedom (CHMM-t), and Laplace (CHMM-L). On 2014-2024
SPY in-sample ($T = 2{,}516$) and 2024-2026 out-of-sample ($T = 572$) data,
all three CHMM variants reach Kolmogorov-Smirnov pass rates above $94\%$
in-sample and $80\%$ out-of-sample, outperforming GARCH(1,1), discrete-HMM,
and recurrent deep-generative baselines across seven distributional and
autocorrelation metrics. A rank-reordering copula extension composes on top
of the per-asset CHMM marginals: across six tickers the Student-t copula
(profile-MLE $\nu^{*} = 6$) reduces off-diagonal correlation MAE to $0.027$
against $0.076$ for a Single Index Model baseline while preserving every
per-asset marginal exactly.
\end{abstract}

%% ACM CCS concepts (2012 taxonomy).
\begin{CCSXML}
<ccs2012>
<concept>
<concept_id>10010147.10010257.10010293.10010294</concept_id>
<concept_desc>Computing methodologies~Neural networks</concept_desc>
<concept_significance>100</concept_significance>
</concept>
<concept>
<concept_id>10003752.10003809.10010172</concept_id>
<concept_desc>Theory of computation~Dynamic programming</concept_desc>
<concept_significance>300</concept_significance>
</concept>
<concept>
<concept_id>10010147.10010257.10010258.10010259</concept_id>
<concept_desc>Computing methodologies~Unsupervised learning</concept_desc>
<concept_significance>500</concept_significance>
</concept>
</ccs2012>
\end{CCSXML}

\ccsdesc[500]{Computing methodologies~Unsupervised learning}
\ccsdesc[300]{Theory of computation~Dynamic programming}
\ccsdesc[100]{Applied computing~Economics}

\keywords{continuous hidden Markov model, Baum-Welch, stylized facts,
volatility clustering, Student-t copula, Single Index Model,
synthetic financial data}

\maketitle

% =================================================================
\section{Introduction}
\label{sec:intro}

Synthetic financial time series are increasingly used for back-testing,
stress testing, and training of downstream risk and trading models. A
generator is useful for these tasks only if it reproduces the three
canonical stylized facts of equity returns: heavy-tailed marginals,
near-zero linear autocorrelation of returns, and persistent positive
autocorrelation of absolute returns. Classical Gaussian hidden Markov
models are widely believed to fail on the third fact because, as
Ryd\'{e}n et al.~\cite{ryden1998stylized} showed at state counts
$K \in \{2, 3\}$, the geometric sojourn-time distribution decays too
fast to sustain volatility clustering.

The recent discrete-HMM framework of Alswaidan et
al.~\cite{alswaidan2026hybrid} works around this limitation by
\emph{quantizing} returns into Laplace quantile bins, frequency-counting
transitions between bins, and adding a Poisson jump mechanism to
reach observed tail heaviness. The quantization makes Kolmogorov-Smirnov
(KS) tests fail by construction on the simulated return series; the
jump mechanism introduces two free hyperparameters ($\epsilon$ for
jump probability, $\lambda$ for mean jump duration).

This paper shows that both workarounds are unnecessary. A continuous
HMM trained by Expectation-Maximization (EM) on raw returns at a single
moderate state count ($K = 18$, selected by a seven-metric sweep) matches
all three stylized facts, passes KS / Anderson-Darling tests uniformly,
and eliminates both hyperparameters of the prior discrete
baseline. Three emission families (Gaussian, Student-t with per-state
$\nu$, Laplace) are compared under identical training; the Student-t and
Laplace variants close the residual kurtosis gap of the Gaussian variant
without sacrificing autocorrelation fidelity.

To extend the framework from single-asset to multi-asset synthesis we
introduce the distinction between \emph{Pipeline A} (per-ticker
independent CHMM fits, no cross-asset coupling) and \emph{Pipeline B}
(SIM~\cite{sharpe1963simplified} and Gaussian / Student-t
copula~\cite{sklar1959fonctions, demarta2005tcopula} constructions on top
of Pipeline A marginals). Pipeline A answers the univariate question
(does the CHMM fit each ticker); Pipeline B answers the joint question
(how well do the marginals compose).

% =================================================================
\section{Methods}
\label{sec:methods}

\subsection{CHMM Definition and EM Updates}
\label{sec:methods_chmm}

Denote the annualized daily excess log-return series by $G_{1:T}$. A CHMM
with $K$ hidden states has parameters $\theta = \{\pi, T, \{\phi_k\}\}$:
initial state distribution $\pi$, $K \times K$ transition matrix $T$ with
$T_{ij} = P(s_{t+1} = j \mid s_t = i)$, and per-state emission
distributions $\phi_k$. Three emission families are considered:
\begin{align}
  \text{CHMM-N}: \quad & G_t \mid s_t = k \sim N(\mu_k, \sigma_k^2), \\
  \text{CHMM-t}: \quad & G_t \mid s_t = k \sim t_{\nu_k}(\mu_k, \sigma_k), \\
  \text{CHMM-L}: \quad & G_t \mid s_t = k \sim \text{Laplace}(\mu_k, b_k).
\end{align}
Training is by standard log-space forward-backward recursions. The
M-step is closed-form Baum-Welch for CHMM-N; Expectation-Conditional
Maximization (ECM) with a golden-section search over $\nu_k \in
[2.5, 100]$ for CHMM-t~\cite{peel2000robust, liu1995ml}; and weighted
median and weighted mean-absolute-deviation estimators for
CHMM-L. All three families use a quantile-based initialization that
places state means on equi-spaced quantiles of $G_{1:T}$. Convergence
tolerance is $10^{-4}$ on the log-likelihood; default iteration count
is 60.

\subsection{Cross-Asset Dependence (Pipeline B)}
\label{sec:methods_ca}

Let $\mathbf{R}_{\text{IS}} \in \mathbb{R}^{T \times d}$ be the $d$-asset
return matrix and let the market-factor column be indexed by $m$. The
SIM construction is
\begin{equation}
  G_{j,t} = \alpha_j + \beta_j G_{m,t} + \eta_{j,t},
  \label{eq:sim}
\end{equation}
fitted by OLS per non-market asset $j \neq m$, with simulation via
$\hat{G}_{j,t}^{(p)} = \hat\alpha_j + \hat\beta_j G_{m,t}^{(p)} +
\hat\eta_{j,t}^{(p)}$ where $G_{m,t}^{(p)}$ is an independent CHMM draw of
the market factor and $\hat\eta_{j,t}^{(p)}$ is an empirical residual
resampled with replacement.

The copula constructions use Sklar's theorem~\cite{sklar1959fonctions}:
any joint CDF with continuous marginals factors as $H = C(F_1, \ldots,
F_d)$. We probit-transform $u_{j,t} = \text{rank}(g_{j,t}) / (T + 1)$,
estimate $\Sigma$ from pairwise Kendall's $\tau$ via $\rho_{ij} =
\sin(\pi \tau_{ij} / 2)$~\cite{embrechts2002correlation}, and for the
Student-t copula select $\nu$ by profile MLE over $\{2, 3, 4, 5, 6, 8,
10, 15, 20, 30\}$. Simulation reorders the ranks of each column of a
per-asset CHMM draw to match the ranks of a copula
sample~\cite{iman1982distribution}, which preserves every asset's
marginal exactly while injecting cross-asset dependence through rank
coupling.

\subsection{Evaluation}
\label{sec:methods_eval}

Every generator produces 1{,}000 independent simulated paths of length
$T$ (200 for Pipeline B, given the quadratic copula cost). Five
complementary metrics score each path before aggregation: KS and
Anderson-Darling pass rates at $\alpha = 0.05$; excess kurtosis;
mean absolute error of the absolute-return autocorrelation function at
lags 1-252 ($\text{ACF-MAE} = \frac{1}{L} \sum_{\tau = 1}^{L} |
\hat\rho^{\text{obs}}_{|G|}(\tau) - \bar{\hat\rho}^{\text{sim}}_{|G|}(\tau) |$);
and Wasserstein-1 distance between pooled simulated and observed
quantiles. Pipeline B additionally reports the Frobenius norm and
off-diagonal MAE of the simulated Pearson correlation matrix relative
to observed.

% =================================================================
\section{Empirical Study}
\label{sec:empirical}

\subsection{Data}

We use daily SPY prices from 2014-01-03 through 2024-01-03 as the
in-sample (IS) window ($T_{\text{IS}} = 2{,}516$) and 2024-01-04 through
2026-04-17 as the out-of-sample (OoS) window ($T_{\text{OoS}} = 572$).
Annualized excess log-returns are computed as $G_t = (1/\Delta t)
\ln(P_t / P_{t-1}) - r_f$ with $\Delta t = 1/252$ and $r_f = 0$. The
observed IS series is leptokurtic (excess kurtosis 7.68), has a linear
ACF of returns that decays to noise within two lags, and an ACF of
$|G_t|$ that remains statistically significant at lag 252.
Figure~\ref{fig:stylized_facts} shows the four canonical stylized-fact
diagnostics on the IS window.

\begin{figure}[t]
\centering
\includegraphics[width=\columnwidth]{sections/figs/Fig-1-Stylized-Facts.pdf}
\Description{Four-panel stylized-facts diagnostic for SPY daily excess
log returns. Panel (a) shows the marginal density is peaked and
heavy-tailed relative to Gaussian and Laplace fits. Panel (b) shows
the normal Q-Q plot with tail deviation on both sides of the
distribution. Panel (c) shows the autocorrelation of returns is near
zero at all lags. Panel (d) shows the autocorrelation of absolute
returns decays slowly and remains positive through lag 252.}
\caption{\textbf{SPY stylized facts} (IS, 2014-01-03 through 2024-01-03,
$T_{\text{IS}} = 2{,}516$). Panel (a): marginal density with Gaussian
and Laplace reference fits. Panel (b): normal Q-Q plot. Panel (c):
ACF of $G_t$. Panel (d): ACF of $|G_t|$. Panels (a-b) evidence heavy
tails; (c) evidences near-zero linear ACF; (d) evidences persistent
volatility clustering.}
\label{fig:stylized_facts}
\end{figure}

\subsection{Model selection: $K = 18$}

A sweep over $K \in \{3, 6, 9, 12, 15, 18, 21\}$ under CHMM-N shows that
$K = 18$ wins or ties on every distributional pass-rate metric (KS IS,
KS OoS, AD IS, AD OoS, Wasserstein-1) and produces the highest simulated
kurtosis (5.11 against observed 7.68). Margins between adjacent $K$
values are under 1 percentage point on KS; the choice is supported by
both the in-sample and out-of-sample pass-rate rankings.

\subsection{SPY model comparison (Pipeline A)}

Table~\ref{tab:spy_comparison} compares CHMM at $K = 18$ against the
discrete-HMM baseline of~\cite{alswaidan2026hybrid} (NJ: no jumps; WJ:
with Poisson jumps, $\epsilon = 0.01$, $\lambda = 3$) and four classical
baselines. The CHMM variants achieve 94.7-95.2\% IS KS pass rate and
80.4-84.0\% OoS KS pass rate, in contrast to the 0\% IS KS of both
discrete variants (a direct consequence of returning bin-centroid
values) and the 23.1\% of GARCH(1,1). The i.i.d.\ bootstrap is the
non-parametric ceiling at 99.6\% IS KS but fails the temporal
diagnostic (ACF-MAE 0.0627, the highest in the panel), confirming
that good marginal fidelity alone does not certify a generator for
time-series downstream use.

\begin{table}[t]
\centering
\small
\caption{SPY model comparison, Pipeline A. 1{,}000 simulated paths,
$\alpha = 0.05$, $K = 18$ for CHMM. IS: 2014-01-03 through 2024-01-03,
$T_{\text{IS}} = 2{,}516$. OoS: 2024-01-04 through 2026-04-17,
$T_{\text{OoS}} = 572$. Observed IS kurtosis 7.68. Column legend:
KS = Kolmogorov-Smirnov pass rate, AD = Anderson-Darling pass rate,
Kurt = mean simulated excess kurtosis,
ACF-MAE = mean absolute error of $\text{ACF}(|G_t|)$ over 252 lags,
$W_1$ = Wasserstein-1 distance (IS).}
\label{tab:spy_comparison}
\begin{tabular}{l cc cc c c c}
\toprule
 & \multicolumn{2}{c}{KS (\%)} & \multicolumn{2}{c}{AD (\%)} & & & \\
\cmidrule(lr){2-3} \cmidrule(lr){4-5}
Model & IS & OoS & IS & OoS & Kurt & ACF-MAE & $W_1$ \\
\midrule
Bootstrap         & 99.6 & 91.0 & 99.4 & 94.2 & 7.43  & 0.0627 & 0.065 \\
Gaussian i.i.d.   & 0.0  & 2.0  & 0.0  & 0.0  & -0.00 & 0.0627 & 0.411 \\
Discrete NJ       & 0.0  & 38.1 & 0.0  & 65.9 & 0.64  & 0.0617 & 0.314 \\
Discrete WJ       & 0.0  & 36.2 & 0.0  & 54.3 & 0.52  & 0.0618 & 0.326 \\
GARCH(1,1)        & 23.1 & 58.3 & 8.5  & 55.5 & 6.97  & 0.0492 & 0.245 \\
CHMM-N            & 94.7 & 84.0 & 86.5 & 76.5 & 5.02  & 0.0513 & 0.110 \\
CHMM-t            & 94.7 & 83.4 & 89.9 & 79.0 & 17.34 & 0.0554 & 0.102 \\
CHMM-L            & 95.2 & 80.4 & 93.2 & 79.9 & 6.76  & 0.0564 & 0.102 \\
\bottomrule
\end{tabular}
\end{table}

\subsection{Six-ticker generalization (Pipeline A)}

Table~\ref{tab:six_ticker} generalizes the CHMM-N result from SPY to
five additional tickers spanning distinct risk profiles: NVDA
(high-volatility semiconductors), JNJ (defensive healthcare), JPM
(financials), AAPL (large-cap tech), QQQ (tech-heavy index). Each
ticker is fit independently with its own CHMM. IS KS pass rates are
uniformly above 95\%. OoS KS pass rates are above 90\% on four of six
tickers; NVDA (55.8\%) and JPM (49.8\%) show OoS degradation
attributable to the 2024-2026 macroeconomic regime (the AI-driven
semiconductor rally for NVDA, the interest-rate-repricing cycle for
financials).

\begin{table}[t]
\centering
\small
\caption{Per-ticker CHMM-N generalization, Pipeline A. $K = 18$,
1{,}000 paths, $\alpha = 0.05$. Each ticker is fit independently; no
cross-asset coupling. Complement: Table~\ref{tab:dependence} evaluates
cross-asset dependence on top of the same marginals.}
\label{tab:six_ticker}
\begin{tabular}{l cc cc c}
\toprule
 & \multicolumn{2}{c}{KS (\%)} & \multicolumn{2}{c}{AD (\%)} & \\
\cmidrule(lr){2-3} \cmidrule(lr){4-5}
Ticker & IS & OoS & IS & OoS & Kurt obs \\
\midrule
SPY    & 95.5 & 81.4 & 88.9 & 76.5 & 7.7 \\
NVDA   & 96.7 & 55.8 & 92.6 & 51.5 & 5.5 \\
JNJ    & 99.4 & 93.9 & 98.8 & 92.1 & 9.0 \\
JPM    & 97.8 & 49.8 & 94.4 & 48.6 & 7.6 \\
AAPL   & 98.3 & 93.8 & 96.9 & 93.5 & 3.2 \\
QQQ    & 95.5 & 92.3 & 87.8 & 95.0 & 3.4 \\
\midrule
Median & 96.7 & 87.1 & 93.5 & 84.3 & 6.6 \\
\bottomrule
\end{tabular}
\end{table}

A quarterly walk-forward refit of the CHMM-N on the most-degraded
ticker (JPM) lifts OoS KS from 49.0\% to 64.3\% and OoS AD from
48.6\% to 63.6\%, a 15-percentage-point recovery, confirming that
the stationarity assumption of the fixed-IS fit is the dominant driver
of the OoS gap on this ticker.

\subsection{Cross-asset dependence (Pipeline B)}

Table~\ref{tab:dependence} evaluates three dependence constructions on
top of the same CHMM-N marginals (Table~\ref{tab:six_ticker}): SIM with
SPY as market factor, Gaussian copula, and Student-t copula with
$\nu^* = 6$ selected by profile MLE. The two copula constructions
preserve every per-asset marginal exactly (mean IS KS
97.5\% / 96.2\%); SIM's affine propagation degrades JPM (31.5\% IS KS)
and AAPL (64.0\%) because the factor-resampled residuals re-enter each
asset's marginal and thicken the tails beyond the asset's own
distribution. On cross-asset correlation reproduction, the Student-t
copula is the best (off-diagonal MAE 0.027, Frobenius 0.184), the
Gaussian copula is second (0.031, 0.206), and the SIM trails (0.076,
0.527). The profile-MLE choice of $\nu^* = 6$ implies non-trivial joint
tail dependence and is in the $4 \leq \nu \leq 10$ range typically
reported for equity-index portfolios~\cite{demarta2005tcopula,
mcneil2015quantitative}.

\begin{table}[t]
\centering
\small
\caption{Pipeline B. Cross-asset dependence on fixed CHMM-N marginals.
200 simulated paths. Rows 1-6: per-asset IS KS (\%), $\alpha = 0.05$
(sanity on the marginals after dependence injection). Rows 7-8:
cross-asset correlation reproduction, $\|\hat\Sigma_{\text{sim}} -
\hat\Sigma_{\text{obs}}\|_F$ and off-diagonal MAE, averaged over paths.
Marginals identical across columns; only the dependence construction
varies.}
\label{tab:dependence}
\begin{tabular}{l c c c}
\toprule
 & SIM & Gaussian & Student-t \\
 &     & copula   & copula ($\nu^* = 6$) \\
\midrule
SPY          & 91.5 & 94.5 & 95.5 \\
NVDA         & 71.5 & 96.0 & 93.0 \\
JNJ          & 88.5 & 100.0 & 100.0 \\
JPM          & 31.5 & 99.5 & 96.5 \\
AAPL         & 64.0 & 99.5 & 99.0 \\
QQQ          & 85.5 & 95.5 & 93.0 \\
\midrule
Mean KS      & 72.1 & 97.5 & 96.2 \\
Median KS    & 78.5 & 97.8 & 96.0 \\
\midrule
Frobenius    & 0.527 & 0.206 & \textbf{0.184} \\
Off-diag MAE & 0.076 & 0.031 & \textbf{0.027} \\
\bottomrule
\end{tabular}
\end{table}

\begin{figure}[t]
\centering
\includegraphics[width=\columnwidth]{sections/figs/Fig-Cross-Asset-Correlation.pdf}
\Description{Four-panel heatmap grid comparing observed and simulated
cross-asset correlation matrices across six tickers. Panel (a) is the
observed IS correlation. Panel (b) is the path-averaged correlation
under the Single Index Model. Panel (c) is under the Gaussian copula.
Panel (d) is under the Student-t copula. The SIM panel visibly differs
from the observed panel in several off-diagonal cells, while both
copula panels are nearly indistinguishable from the observed panel.}
\caption{\textbf{Pipeline B cross-asset correlation reproduction}
(companion to Table~\ref{tab:dependence}; IS window,
$T_{\text{IS}} = 2{,}516$, six tickers SPY, NVDA, JNJ, JPM, AAPL, QQQ).
Panel (a): observed correlation matrix on IS returns. Panel (b):
path-averaged correlation under the SIM with SPY as market factor.
Panel (c): path-averaged correlation under the Gaussian copula on
CHMM-N marginals. Panel (d): path-averaged correlation under the
Student-t copula with $\nu^* = 6$. Each simulated panel averages 200
paths of length $T_{\text{IS}}$. Color scale in $[-1, 1]$, red-white-blue
diverging; closer visual match to panel (a) is better. The SIM panel
shows visibly larger residuals off the diagonal, matching the
numerical ordering.}
\label{fig:cross_asset_corr}
\end{figure}

% =================================================================
\section{Discussion and Conclusion}
\label{sec:discussion}

At a single moderate state count ($K = 18$) a continuous HMM trained by
standard EM reproduces all three canonical stylized facts of equity
returns without jump augmentation, without return quantization, and
without additional tunable hyperparameters. Against the discrete-HMM
baseline of~\cite{alswaidan2026hybrid} the IS KS pass-rate gap is 94.7
percentage points and is attributable to the move from bin-centroid
emissions to continuous emissions: continuous emissions are a
prerequisite for the simulated series to pass KS at all, and the
elimination of Poisson-jump augmentation removes two
hyperparameters. The gap between Ryd\'{e}n et al.~\cite{ryden1998stylized}
(K = 2-3, failure to reproduce absolute-return ACF) and the present
result (K = 18, ACF-MAE 0.0513) is a direct consequence of state
resolution: sufficient diagonal dominance across a larger number of
states approximates the slow ACF decay through a sum-of-exponentials,
whereas K = 2-3 cannot.

Pipeline B's ordering (Student-t copula over Gaussian copula over SIM)
replicates the ordering reported by~\cite{alswaidan2026hybrid} with
discrete marginals, suggesting the rank-reordering advantage is
intrinsic to the construction and transfers across marginal families.
Stationarity is the principal limitation: OoS pass rates drop on NVDA
and JPM under the 2024-2026 regime, and a quarterly walk-forward refit
recovers roughly 15 percentage points on JPM, supporting stationarity
as the mechanism rather than mis-specified state count or
emission family.

\paragraph{Scope.} Option pricing and implied-volatility-surface
calibration are out of scope. The six-asset universe is sufficient to
establish the SIM-versus-copula ordering but does not address the
scaling behavior of the copula construction to hundreds of assets,
which would likely require vine or factor copulas~\cite{bedford2002vines}.

\paragraph{Ethics and data provenance.} All data used in this paper
are daily OHLCV (open, high, low, close, volume) closing prices and
volume-weighted average prices for six publicly-traded U.S.\ equities
(SPY, NVDA, JNJ, JPM, AAPL, QQQ) retrieved from a commercial market-data
vendor under a standard data-use license. No personally-identifying or
privacy-sensitive information is used. The study presents no human-subject
research.

%% ACM-style bibliography.
\bibliographystyle{ACM-Reference-Format}
\bibliography{References_v8-icaif}

\end{document}
